\chapter{A Human-Centric Model for Understanding LLM Errors}
\label{sec:ch13-error-model}

\section*{Setting the Scene}
\addcontentsline{toc}{section}{Setting the Scene}

\paragraph{The problem.}
The thirteen chapters of this book have given you the vocabulary of
the engineer who builds, trains, deploys, and governs a Large
Language Model. That vocabulary contains terms like
\emph{cross-entropy}, \emph{prompt injection}, \emph{lost in the
middle}, \emph{reward hacking}, and \emph{retrieval grounding}, and
each of those terms picks out a real and important failure mode.
What that vocabulary does not give you, on its own, is a way of
speaking to the user across the desk who has just been handed a
wrong answer by a language model and wants to know, in plain
language, why. ``The cross-entropy objective trains for plausibility
rather than truth'' is correct, but it is not what the user came to
hear. The problem of this chapter is the gap between the engineer's
vocabulary, which is precise, and the user's experience of error,
which is unmediated. We want a model of LLM failure that an
intelligent non-specialist can hold in their head, that maps onto
the technical mechanisms covered in previous chapters, and that
gives a user something to do when an answer arrives that they
suspect is wrong.

\paragraph{How we got here.}
The study of language as a layered phenomenon is older than the
language model. Linguists since at least the early twentieth
century have distinguished between the structure of an utterance
(its syntax), the meaning it conveys (its semantics), and the
situational effect it has (its pragmatics). Each level has its own
rules, its own characteristic failures, and its own diagnostic
vocabulary. Noam Chomsky's distinction between competence and
performance, in the 1960s, separated the rules of grammar from the
noisy reality of speech \citep{chomsky1965aspects}. H.~Paul Grice,
in his 1967 William James lectures, articulated the cooperative
principle and the conversational implicatures that govern what a
speaker means as opposed to what they literally say
\citep{grice1975logic}. John Searle's speech-act theory categorized
what utterances do as well as what they say
\citep{searle1969speech}. None of these thinkers was talking about
machines, and yet their categories are the ones that turn out to
matter when the machine begins to talk.

The translation to artificial systems began with the question of
whether a symbol-manipulating system could ever truly mean anything
at all. Stevan Harnad's 1990 paper \emph{The Symbol Grounding
Problem} \citep{harnad1990symbol} asked how the symbols inside a
computer could come to refer to things in the world, given that the
symbols themselves are just shapes. The question became urgent when
statistical language models began producing fluent text without any
obvious mechanism for connecting that text to anything outside the
corpus. Emily Bender and Alexander Koller's 2020 paper
\emph{Climbing towards NLU} \citep{bender2020climbing} sharpened the
argument: a system trained only on linguistic form cannot, by that
training alone, learn meaning, because meaning is a relation
between form and the world that the form alone does not contain.
Gary Marcus, in a series of papers and books
\citep{marcus2020next}, argued that the failures of modern LLMs are
not bugs but the expected consequence of an architecture that has
no concept of reference at all.

The practitioner literature has its own vocabulary for the same
problem. \emph{Hallucination} names the case in which a model
produces fluent text that has no relationship to fact
\citep{ji2023survey}. \emph{Prompt brittleness} names the case in
which two prompts a human would read as equivalent produce
qualitatively different outputs. \emph{Lost in the middle} names
the case in which a model neglects information sitting in the
middle of a long context. These terms are useful for engineers but
they are jargon, and they do not, considered as a set, form a
coherent taxonomy. They are a list of phenomena rather than a model
of the underlying structure.

\paragraph{Where we stand.}
The model this chapter develops has five axes, arranged from the
mechanical to the human. The first is \emph{structure} --- whether
the output is well-formed for its purpose. The second is
\emph{meaning} --- whether the output is internally coherent and
says what it appears to say. The third is \emph{context} ---
whether the output is appropriate for the situation, the purpose,
and the audience. The fourth is \emph{grounding} --- whether the
output connects to verifiable reality or has drifted into
plausible-sounding fabrication. The fifth is \emph{environment} ---
whether the prompt itself contained enough of what the human in the
situation knows for the question to be answerable at all. The five
axes form a ladder. At the bottom, the model is at its strongest
and failures are rare. At the top, the model is, in a precise
sense, at the mercy of the human, and failures are inevitable
unless the human compensates.

The pedagogical use of the model is twofold. For users, it is a
diagnostic ladder: when an answer arrives that smells wrong, walk
the axes and ask which one has failed. For engineers, it is a way
of mapping the vocabulary of earlier chapters onto a shape that
non-specialists can hold. Cross-entropy training is a structure-
and-meaning mechanism that makes no native attempt at grounding;
retrieval is a grounding mechanism that does not by itself solve
context; tool use is an environment-importing mechanism whose
contract is structural. The ladder makes these connections
explicit.

\paragraph{What goes wrong.}
The model itself has failure modes worth disclosing in advance.
The first is the temptation to use it mechanically, as if the
diagnosis of an output were a checklist that always returns a
single answer. Real failures often span multiple axes: a model
that is wrong about a fact and presents it in the wrong tone has
failed on grounding and context simultaneously, and pulling the
two apart can be artificial. The second is the temptation to treat
the axes as orthogonal. They are not. Failures of structure tend to
drag meaning with them; failures of context tend to drag grounding.
The third is the deeper mistake of treating the model as a solution
rather than as a diagnostic vocabulary. The model does not, by
itself, fix anything. It tells you which earlier chapter to
consult.

\section{Learning Objectives}

By the end of this chapter, you should be able to:
\begin{itemize}
  \item Articulate the five axes of the human-centric error model
        (structure, meaning, context, grounding, environment) and
        explain why they are arranged in the order they are.
  \item Diagnose a wrong LLM output by walking the ladder and
        identifying which axis has failed.
  \item Connect each axis to the underlying technical mechanism
        covered in earlier chapters of this book.
  \item State the limits of the model: where the axes overlap,
        where diagnosis is interpretive rather than algorithmic,
        and what kinds of failure escape the framework entirely.
  \item Use the model to communicate with non-specialist
        stakeholders --- users, regulators, auditors, executives
        --- about LLM errors in language they can act on.
\end{itemize}

\section{LLM Components Covered}

\textbf{Diagnostic framework}
\begin{itemize}
  \item Five-axis ladder: structure, meaning, context, grounding,
        environment.
  \item Mapping from each axis to earlier-chapter mechanisms.
  \item Walking the ladder as a structured diagnostic procedure.
\end{itemize}

\textbf{Communication relevance}
\begin{itemize}
  \item Translation between the engineer's vocabulary and the
        user's experience of error.
  \item Failure-class language for incident reports, audits, and
        stakeholder conversations.
\end{itemize}

\textbf{Systems relevance}
\begin{itemize}
  \item Locus of intervention: which axes call for model-side
        change versus system-side change.
  \item Limits of the framework: normative, adversarial, and
        drift failures that escape the five axes.
\end{itemize}

\section{Notation and Standing Assumptions}
\label{sec:ch13-notation}

\begin{notationbox}
\textbf{Symbols.}
Let $y$ denote a model's output, $p$ the prompt the model sees,
$s$ the user's actual situation and intent, $w$ the state of the
world (\emph{ground truth}), and $h$ the human's full mental
context (everything they know about the situation, including
things not in $p$). The five axes can be expressed as predicates
over these objects:
\begin{align}
  \mathrm{structure}(y) &= \mathbb{1}\!\left[y \in \mathcal{V}_{\mathrm{form}}\right] \label{eq:ch13-structure}\\
  \mathrm{meaning}(y) &= \mathbb{1}\!\left[y \text{ is internally consistent}\right] \label{eq:ch13-meaning}\\
  \mathrm{context}(y; s) &= \mathbb{1}\!\left[y \text{ is appropriate for } s\right] \label{eq:ch13-context}\\
  \mathrm{grounding}(y; w) &= \mathbb{1}\!\left[y \text{ is true in } w\right] \label{eq:ch13-grounding}\\
  \mathrm{environment}(p; h) &= \mathbb{1}\!\left[p \text{ encodes enough of } h\right] \label{eq:ch13-environment}
\end{align}
where $\mathcal{V}_{\mathrm{form}}$ is the set of well-formed
outputs for the purpose at hand (a JSON schema, a grammatical
sentence, a correct LaTeX file, and so on). The first four
predicates are evaluated on the model's output. The fifth is
evaluated on the prompt --- it asks whether the human supplied the
model with enough of their environment for the question to be
answerable. An output that fails on any axis below environment is
a candidate model-side fix; an output that fails on environment is
a candidate prompt-side or system-side fix.

\textbf{Standing assumptions.}
\textbf{(A12.1)} The five axes are not orthogonal; real failures
frequently span multiple axes. The taxonomy is a diagnostic
vocabulary, not an algebraic decomposition.
\textbf{(A12.2)} The ladder ordering (structure $\to$ meaning $\to$
context $\to$ grounding $\to$ environment) reflects the empirical
tendency of modern LLMs: they are strong at the bottom and weak at
the top. The ordering is a heuristic, not a theorem.
\textbf{(A12.3)} Training alone is \emph{insufficient} to close
the grounding and environment axes: the model cannot, by training
on additional text, verify a claim against the present-day world
or recover context the user has not written into the prompt.
Inference-time mechanisms can close part of the gap (retrieval,
tool use, agentic critique loops, calibrated abstention,
clarifying-question training), and the chapter assumes these are
system-level interventions (Chapter~10, Chapter~12) layered on top
of the model rather than properties internal to it.
\end{notationbox}

\section{The Five Axes}

We treat each axis in turn. For each, we give a definition, an
illustrative failure, the technical mechanism in earlier chapters
that produces or remedies the failure, and a note on what diagnosis
along that axis looks like in practice.

\subsection{Structure: is the output well-formed?}

A model's output is structurally correct if it satisfies the formal
constraints of the format being requested. A JSON object that
parses; a grammatically complete sentence; a Python script that
imports without error; a SQL query that the database accepts; a
LaTeX document that compiles. Structure is the most mechanical of
the five axes and is, for that reason, where modern LLMs perform
best.

A structural failure is the kind a parser can catch. The model
produces a JSON-like blob with a missing comma. The model writes
Python whose indentation contradicts the syntax rules. The model
opens a markdown code fence and forgets to close it. These failures
are real but, in the current generation of frontier models,
increasingly rare. They become more common under three kinds of
pressure: very long outputs, formats with complex nested structure
(deeply nested JSON, Lisp, formal grammars), and unfamiliar formats
outside the training distribution.

The technical mechanism is, broadly speaking, the
tokenization-and-decoding pipeline of Chapters~3 and~5. Modern
serving stacks reduce the structural failure rate further by
\emph{constrained decoding}: the sampler is restricted at each step
to tokens that keep the partial output within the target grammar.
JSON-mode and JSON-schema enforcement, regex-constrained generation,
and finite-state-machine guards over the token distribution are all
instances of the same idea. Function-calling APIs (Chapter~12)
typically sit on top of constrained decoders.

In diagnosis, structural failures are the easiest to identify ---
they are the failures that automated tooling already catches. They
are also the easiest to be falsely reassured by: a syntactically
valid JSON object can still mean the wrong thing, say the wrong
thing, and refer to a thing that does not exist. A practitioner
who confirms structure and stops has confirmed the least
interesting property of the output.

\subsection{Meaning: is the output internally coherent?}

A structurally correct output can still fail at the level of
meaning. The model produces a paragraph whose first sentence
contradicts its third. The model defines a variable in its code
and then references a different variable name in the next line.
The model's argument relies on a premise it has elsewhere denied.
These failures are detectable by careful reading without any
appeal to the world; they are failures of internal coherence.

The mechanism that produces meaning failures is the same mechanism
that produces fluent text in the first place. A language model
trained on cross-entropy (Chapter~1) optimizes for the
\emph{plausibility} of the next token given the prefix.
Plausibility is a local property. A token can be plausible at
every position in a sequence and still produce, at the end of the
sequence, a global contradiction. A single forward pass has no
architectural verifier; the model emits each token conditioned on
the prefix without checking the output as a whole against itself,
and coherence is a side-effect of training rather than an enforced
invariant.

Recent training narrows the gap. Reasoning-trained models --- the
OpenAI o1/o3 line, DeepSeek R1, Claude with extended thinking ---
learn to walk back contradictions in their own reasoning traces
during generation, in effect performing a form of self-checking
inside the inference call. Outside the model, inference-time
techniques sample multiple completions and compare them
(self-consistency), invoke the model to critique and revise its
own output (self-refine, Reflexion), or hand the output to a
separate verifier model trained for this purpose. Agentic systems
(Chapter~12) extend the same pattern by adding explicit
critique-and-revise loops to a chain of LM calls. None of these is
a guarantee in the way a JSON parser is for structure; each is a
mitigation that trades inference compute and orchestration
complexity for reduced inconsistency.

This is the axis that benchmarks tend to overstate. A model that
scores 85 on a multiple-choice reading-comprehension test may
nevertheless produce, in free generation, paragraphs that
contradict themselves in subtle ways. Domain expertise tends to
unmask this: the technical reader who knows the territory catches
inconsistencies that the casual reader does not. The
\emph{lexical-semantic-drift} literature documents the same
phenomenon at a smaller granularity --- the same term used with
slightly different meanings across a long output.

In diagnosis, the question to ask is not whether the output sounds
right but whether it \emph{is} right considered as a self-contained
artifact. Read it once for fluency, then again for contradictions.
When the cost of a meaning failure justifies the additional
compute, engage the appropriate inference-time mitigation from the
list above. These are, in the language of Chapter~6, ways of
trading inference compute against meaning reliability.

\subsection{Context: is the output appropriate for the situation?}

The third axis is whether the output is appropriate for its
situation, purpose, and audience. A model can produce a
structurally valid, internally coherent paragraph that nevertheless
answers the wrong question, addresses the wrong audience, strikes
the wrong tone, or pitches the wrong level of detail. The
technical literature calls this \emph{pragmatic competence}; the
user recognizes it as the model having ``answered a different
question.''

Examples are easy to find. A user asks for a summary of a research
paper for a non-specialist audience and receives a paragraph dense
with jargon. A user asks a quick yes-or-no question and receives a
five-paragraph essay with a yes-or-no buried in the middle. A user
asks a clarifying question about a sensitive topic and receives a
generic safety disclaimer instead of the requested clarification.
None of these are factually wrong. All of them are wrong for the
situation.

The mechanism is largely the gap between what the user knows about
the situation and what the prompt encodes. The model has access
only to the tokens in its context window. Anything the user knows
about themselves, their audience, their goals, or their constraints
that is not written into the prompt does not exist for the model.
Few-shot prompting (Chapter~6) is, among other things, an attempt
to encode situational expectations through examples. System prompts
(Chapter~12) are an attempt to fix the situation up front.
Alignment training (Chapter~9) shapes the model's defaults along
this axis --- politeness, refusal, length, formality --- but the
defaults are defaults, not the situation.

In diagnosis, context failures are easy to recognize and harder to
fix. The user can usually point at the answer and say what is
wrong with it (``too long,'' ``too technical,'' ``misses the
actual question''). The fix is almost always to give the model
more of the situation: who you are, who you are writing for, why
you need the answer, what would count as a good one. The asymmetry
between the ease of recognition and the difficulty of fix is the
diagnostic signature of a context failure.

\subsection{Grounding: does the output connect to reality?}

The fourth axis is whether the output connects to verifiable
reality or has drifted into plausible-sounding fabrication. The
practitioner literature calls a grounding failure a
\emph{hallucination} \citep{ji2023survey}: the model produces a
fluent, confident, well-formed claim that does not correspond to
any fact in the world.

Grounding failures are uniquely dangerous because no internal
reading of the output reveals them. A hallucinated citation has
the same structural form as a real one. A hallucinated statistic
has the same surface plausibility as a true one. A hallucinated
historical fact reads with the same confidence as a verified one.
The check must come from outside the text. The model has produced
the most probable continuation of its prompt; the most probable
continuation is sometimes true, sometimes false, and the model
cannot, on its own, tell the two apart.

The mechanism is the cross-entropy objective itself. A model
trained to predict the next token given the prefix optimizes for
textual plausibility, not for truth. There is no signal in the
training loss that distinguishes a fact the model has correctly
memorized from a confabulation that fits the surrounding pattern.
The deeper problem, articulated by \citet{harnad1990symbol} and
sharpened by \citet{bender2020climbing}, is that a system trained
only on form has no native way to ground form in reference. The
literature distinguishes \emph{intrinsic} hallucination, in which
the output contradicts source material the model was given, from
\emph{extrinsic} hallucination, in which the output introduces
unverifiable claims. Both are grounding failures.

The remedy at the system level is retrieval (Chapter~10). A
retrieval-augmented generation pipeline replaces the question
``what is the most probable thing to say'' with ``what is the most
probable thing to say given these specific sources.'' The grounding
is no longer the model's problem; it is the retriever's problem
and the prompt's problem. Tool use (Chapters~9 and~11) extends the
same idea to verification: a model that can call a calculator can
delegate arithmetic, a model that can query a database can
delegate fact lookup. The pattern in each case is to take the work
that demands grounding and move it outside the language model
entirely.

In diagnosis, grounding failures call for a particular discipline.
Treat any factual claim by the model as a hypothesis. If the claim
matters, verify it. If it cannot be verified, do not use it. The
careful practitioner builds workflows that make verification cheap
(retrieval, citations, linked sources) so that grounding failures
are catchable by construction.

\subsection{Environment: did the prompt encode the situation?}

The fifth axis sits at the human end of the ladder, and is the
hardest to see because, by construction, it is invisible. An
\emph{environment} failure occurs when the answer is wrong because
the prompt did not contain enough of the situation for the
question to be answerable. The user knows things --- about
themselves, their domain, their constraints, the local conventions
of their organization, the events of yesterday afternoon --- that
the model does not. If the user does not write those things into
the prompt, the model cannot use them. The model can produce a
fluent, internally coherent, grounded, well-formed answer that is
wrong because the question itself was underspecified.

This is the failure mode that experienced LLM users learn to
recognize first and explain to others last. It looks, from the
inside, like a model failure: the answer is wrong, so the model
must be wrong. From the outside, it is a prompt failure: the
model gave the most reasonable answer to the question it was
actually asked, which differs from the question the user thought
they were asking. The gap between the two is the user's
environment.

The mechanism is fundamental and not fixable by training. A
language model has access to its weights and its context window.
It does not have access to the user's calendar, the user's team,
the user's organization's coding conventions, the user's medical
history, the user's project's constraints, or any of the several
hundred other things that shape what counts as a good answer in a
real situation. Some of these can be imported into the context:
through retrieval (Chapter~10), through tool use, through careful
prompting, through long-running agent contexts that accumulate
state. None of them can be imported automatically.

In diagnosis, an environment failure has a characteristic shape.
The user reads the answer and feels frustrated; the answer is, on
its own terms, fine; what is wrong is that it does not address
the real situation. The fix is rarely to ask the model again. The
fix is to write down, explicitly, the parts of the environment the
model needed to see. Practitioners who become good at LLMs are
practitioners who have learned which parts of their environment
need importing for which kinds of question.

\section{The Diagnostic Ladder}

The five axes form a ladder because they are arranged in roughly
increasing distance from what the LLM can reliably do on its own.
The ordering is empirical, not theoretical, but it is robust enough
to use as a default diagnostic procedure.

When an output arrives that strikes the user as wrong, the
recommended walk is bottom-up:
\begin{enumerate}
  \item \emph{Structure}. Does the output parse, compile, validate
        against its schema? If not, the failure is structural and
        the remedy lives in Chapter~6 (decoding) or Chapter~12
        (constrained generation).
  \item \emph{Meaning}. Read the output as a self-contained
        artifact. Does it contradict itself? Does it use a term
        differently in different sentences? Does its argument rely
        on a premise it has elsewhere denied? If yes, the failure
        is one of meaning, and the remedies are sampling multiple
        completions, asking the model to check its own work, or
        using chain-of-thought to externalize the reasoning.
  \item \emph{Context}. Is the output appropriate for the
        situation --- the right tone, the right level, the right
        format, the right question? If not, the failure is
        contextual and the remedy is to import more of the
        situation into the prompt (system prompts, examples,
        explicit role assignment).
  \item \emph{Grounding}. Are the factual claims in the output
        true? If verification is hard, treat the claims as
        hypotheses. If verification is easy, do it. Failures here
        call for retrieval, tool use, or human review.
  \item \emph{Environment}. Is the question itself fully specified?
        Did the prompt contain enough of the user's situation for
        the question to have a single right answer? If not, no
        amount of model improvement will help; the user must
        supply more of their environment.
\end{enumerate}

The discipline of walking the ladder bottom-up is that lower-axis
failures are typically easier to diagnose and cheaper to remedy
than higher-axis ones. A practitioner who jumps straight to ``the
model hallucinated'' may have missed a meaning or context failure
that does not require any change to the model at all. Conversely,
a practitioner who never considers grounding will deploy systems
that are confidently wrong about the world.

When a single output fails on more than one axis, the procedure
is to label all axes that fail (the failure label is a subset
$F \subseteq \{$structure, meaning, context, groundedness,
environment$\}$, not a single name) and to address them in
ladder order: the lowest-axis failure in $F$ is fixed first
because higher-axis fixes assume the lower axes hold. A useful
heuristic for prioritising the work is to ask which axis, if it
had been correct, would have made the others correctable: a
grounded answer in the wrong tone can be fixed by the user, but
a fluent answer about a fact that does not exist cannot. The
limiting axis is the one to pay first.

\section{Connections to Earlier Chapters}

The five axes connect to the technical mechanisms of earlier
chapters in ways that are worth making explicit. Engineers who
hold the ladder in mind have a way of mapping user-facing
complaints onto technical interventions.

\textbf{Structure} is the territory of decoding (Chapter~6),
constrained sampling, and the structured-output contracts of tool
use (Chapter~12). The mathematical objects involved are the
softmax distribution, the JSON-schema-validating sampler, and the
function-calling JSON contract.

\textbf{Meaning} is the territory of cross-entropy training
(Chapter~1), embeddings (Chapter~4), and the inference-compute
trade-offs of self-checking and chain-of-thought (Chapter~6). The
mathematical objects involved are the cross-entropy gradient, the
embedding-space cosine similarity, and the sampling distributions
that drive temperature and majority voting.

\textbf{Context} is the territory of in-context learning
(Chapter~6), prompt construction (Chapter~10), system prompts
(Chapter~12), and the tone-and-style components of alignment
(Chapter~9). The mathematical objects involved are the
preference-model log-likelihood, the few-shot prompt format, and
the RLHF KL penalty that keeps the policy from over-tuning.

\textbf{Grounding} is the territory of retrieval (Chapter~10), tool
use, the symbol-grounding literature \citep{harnad1990symbol,
bender2020climbing}, and the memorization-versus-hallucination
distinction of Chapter~6. The mathematical objects involved are
BM25, the dense-retrieval inner product, the ReAct loop, and the
verification predicates of Chapter~11's hazard analysis.

\textbf{Environment} is the territory of system design rather than
model design. Long-context architectures (Chapter~5) extend how
much environment can be imported. Agent loops (Chapter~12)
accumulate environment over time. Documentation, logging, and
incident reporting (Chapter~11) make the environment the
\emph{system} sees explicit and auditable.

The general pattern is that lower axes are increasingly
\emph{model-side}, while higher axes are increasingly
\emph{system-side}. A team that wants to improve structure will
work on decoding; a team that wants to improve grounding will
work on retrieval; a team that wants to improve environment will
work on their prompt-construction discipline. The ladder makes the
locus of intervention explicit.

\section{What the Model Does Not Capture}

The taxonomy is useful because it is simple. It is also incomplete.
A few classes of failure deserve mention as cases where the model
helps less than one might hope.

The first is the \emph{normative} failure: an output that is
structurally fine, internally coherent, situationally appropriate,
factually grounded, and environmentally informed but that is, by
some standard the user holds, simply wrong --- offensive, biased,
unethical, in conflict with values the user expects the system to
hold. These failures live in Chapter~9's territory rather than
this chapter's. The five axes ask whether the answer is correct;
the normative dimension asks whether the answer is right, and
right is not the same thing as correct.

The second is the \emph{adversarial} failure: an output that has
been deliberately induced by a malicious prompt, a poisoned
retrieval source, or a manipulated tool. These failures cross the
trust boundaries of Chapter~10 and Chapter~12 and are addressed by
defenses in depth rather than by diagnostic ladders.

The third is the \emph{drift} failure: an output that was correct
yesterday and is wrong today because the world has changed. This
is a failure of the system's integration with reality, not of any
single output. The remedy lives in monitoring, evaluation, and
governance (Chapter~11).

The fourth is the \emph{automation complacency} failure: an output
that a human reviewer would have caught if they were producing the
output themselves, but did not catch because they were in a
verification role rather than a production role. This failure
interacts with all five axes and deserves more attention than it
typically receives.

The human-factors literature has documented automation complacency
since the 1980s. Parasuraman and Manzey's review
\citep{parasuraman2010complacency} established the consistent
finding: humans monitoring automated systems miss anomalies at
significantly higher rates than humans performing the same task
manually, with the effect size depending on how plausible the
automated output appears and how rarely errors occur. Both
conditions are routinely present in LLM deployments.

LLM outputs are fluent by construction. A wrong answer arrives in
the same register, the same sentence structure, and the same
apparent confidence as a correct one. The five-axis taxonomy
exists precisely because these failures are hard to detect by
reading — the output must be interrogated, not just consumed.
Under time pressure, a reviewer who sees a fluent, well-structured
response to a familiar question will process it more shallowly
than a reviewer producing the answer from scratch, because shallow
processing is the default mode for text that reads like it is
probably correct. This is not a failure of character or
attention; it is the normal operation of a cognitive system
designed to allocate effort efficiently.

The practical consequence is that human verification of LLM
output is not a reliable substitute for human production of the
same output, even when the reviewer is competent and the task
is familiar. Three patterns make the effect worse in practice:

\begin{itemize}
  \item \emph{Low base rate of errors.} When the LLM is accurate
        95\% of the time, reviewers experience 95 correct outputs
        for every 5 wrong ones. After a sufficient run of correct
        outputs, the prior that the next output is also correct
        becomes strong, and vigilance falls. The errors that matter
        most --- the 5\% --- arrive at the moment of lowest
        vigilance.
  \item \emph{Plausibility of the error class.} The diagnostic
        ladder makes clear that grounding failures --- the model
        asserting a false fact fluently --- are the hardest axis
        to detect. These are also the failures most vulnerable to
        complacency, because a grounding failure is invisible
        unless the reviewer independently knows the fact in
        question.
  \item \emph{Throughput pressure.} Automation is often deployed
        to increase throughput. Higher throughput means less time
        per output for the reviewer. The two effects --- less time
        available, more outputs per hour --- compound the
        complacency effect rather than cancelling it.
\end{itemize}

\Cons\ The engineering implication is that a human-in-the-loop
review step should not be counted as a reliable error-catching
mechanism unless the review conditions are specifically designed
to support active verification: enough time per output to
interrogate rather than scan, periodic calibration cases
(known-error outputs inserted to keep reviewers in an active
detection mode), and metrics on review accuracy, not just review
throughput. Chapter~10 (\S\ref{sec:verification-cost-erosion})
derives the economic conditions under which verification overhead
erodes the benefit of LLM deployment; the complacency effect
raises the effective error rate seen by users above the raw model
error rate, which shifts the break-even threshold further.

The five-axis model is meant for the user-facing case in which a
single output is in front of you and you want to know what to do.
For everything else, the earlier chapters remain the authoritative
reference.

\section{Worked Example: Diagnosing a Wrong Code Suggestion}

Suppose a user asks an LLM coding assistant to write a function
that filters a list of customer records by region and returns the
names sorted by purchase total. The model returns a function. The
user runs it and the output is wrong. We walk the ladder.

\textbf{Structure.} Does the function parse and import? Suppose
yes. Structure passes.

\textbf{Meaning.} Read the function. Does it define what it claims
to define? Does it contradict itself --- declaring a return type
of \texttt{list[str]} but returning tuples? Suppose the function
is internally coherent. Meaning passes.

\textbf{Context.} Is the function appropriate for the situation?
Suppose the user works in a Python codebase that uses a particular
ORM and the model wrote raw list-comprehension code instead. That
is a context failure: the function works in a vacuum but does not
fit the codebase. The remedy is to add the codebase conventions
to the prompt.

Suppose instead the function uses the right ORM. Context passes.

\textbf{Grounding.} Are the references in the code real? If the
function calls \texttt{customers.filter\_by\_region(r)}, does that
method exist on the \texttt{customers} model in the user's actual
ORM? If the model invented a method that looked plausible but does
not exist in the codebase, that is a grounding failure. The remedy
is to give the model the actual ORM definitions, by retrieval or
by citation.

Suppose all the methods are real. Grounding passes.

\textbf{Environment.} Did the prompt say what ``region'' meant in
this codebase? In some businesses, region is a country code; in
others, it is a sales territory id; in still others, it is a free-
text label. If the model assumed the wrong one, the function will
be syntactically and semantically correct, will use the right ORM,
will call real methods, and will still be wrong. That is an
environment failure, and no fix to the model will repair it. The
user must specify what \emph{region} means in their environment.

The walk shows the discipline: we did not stop at the first axis
that passed; we did not jump to the most exotic explanation; we
located the failure at the level where the remedy lives.

It is worth pausing on the grounding failure case to name the
mechanism that produced it, because the ladder's value to an
engineer is greatest when it points back to the chapter where the
fix lives.

The invented ORM method --- \texttt{customers.filter\_by\_region()}
--- is a grounding failure produced by the mechanism of Chapter~1:
next-token prediction trains for plausibility, not for truth.
Method call patterns of the form \texttt{object.verb\_by\_field()}
are common in the training corpus; the model produces what is
statistically likely given the surrounding code structure.
There is no signal in the pretraining objective that distinguishes
a method call that is real in the user's specific ORM version from
one that merely fits the pattern. Notice also \emph{where} the
failure concentrates: not in the high-level structure (filter a
list, sort by numeric field, return names) but in a specific
peripheral detail (the exact method name). This distribution is
characteristic of grounding failures. The skeleton of a correct
answer is more strongly represented in training than the specific,
deployment-relevant details that make the skeleton executable.
Hallucinated citations, invented API endpoints, and fabricated
function signatures all follow the same shape: the genre is correct,
the specific referent is not.

The system-level remedy --- giving the model the actual ORM schema
before generation --- works by changing the question from ``what is
statistically likely'' (a pretraining question) to ``what is
consistent with this document'' (a grounding question). Retrieval
is not a fix to the model; it is a reframing of the task so that
the model's native strength (conditional plausibility) is applied
to a context that already contains the ground truth.

The environment failure case carries a different lesson. No change
to the model, no retrieval pipeline, and no schema injection will
help if the user's intent --- what ``region'' means in their
business --- was never written into the prompt. This failure is
invisible from inside the model and invisible from inside the
system. It is visible only to the user, who knows what they meant.
The diagnostic value of the environment axis is precisely that it
names the class of failure no engineering intervention can prevent:
the failure that belongs to the human end of the loop.

\section{Bridges to Other Weeks}

\begin{itemize}
  \item \textbf{Chapter~1 (Statistical Foundations).} The
        \emph{Meaning} axis is the user-visible consequence of the
        cross-entropy training objective: the model learns
        plausibility, not truth, and global coherence is a
        side-effect rather than an enforced property.
  \item \textbf{Chapter~4 (Tokenization and Embeddings).} The
        embedding space is the substrate on which meaning is
        encoded; the structure axis sits on top of the
        tokenization pipeline.
  \item \textbf{Chapter~5 (Transformer Architecture).} Context
        length and lost-in-the-middle effects are direct
        contributors to the \emph{Environment} axis: the model can
        only use what it can attend to.
  \item \textbf{Chapter~6 (Pretraining and In-Context Learning).}
        In-context learning is the primary mechanism for shaping
        the \emph{Context} axis at inference time; sampling
        strategies trade compute for \emph{Meaning} reliability.
  \item \textbf{Chapter~9 (Alignment).} Alignment shapes the
        defaults along the \emph{Context} axis (tone, refusal,
        length) without changing the \emph{Grounding} axis at all
        --- a confusion this chapter helps disentangle.
  \item \textbf{Chapter~10 (RAG, Tools, Evaluation).} Retrieval is
        the primary system-level remedy for \emph{Grounding};
        tools are the primary remedy for \emph{Environment} when
        external state is needed.
  \item \textbf{Chapter~11 (Governance).} The diagnostic ladder
        feeds incident reports, hazard analyses, and assurance
        cases. A failure that climbs the ladder to
        \emph{Grounding} or \emph{Environment} typically requires
        a system-level control rather than a model-level fix.
  \item \textbf{Chapter~12 (APIs, Agents, MCP).} Agents accumulate
        environment over time; the structured-I/O contracts of
        MCP bound the \emph{Structure} axis tightly enough that
        higher-axis failures dominate the production failure mix.
\end{itemize}

\section{Chapter 13 Takeaway}

The five-axis model is not a theory. It is a vocabulary for
talking about LLM failures across the gap between the engineer's
view and the user's experience. Its merit is that it is short
enough to memorize, ordered enough to use as a procedure, and
connected enough to let a non-specialist locate a failure on the
map. Its limit is that the axes are not orthogonal, the ordering
is heuristic, and some failures escape the framework entirely. Use
it as a diagnostic ladder, not as a proof system. When it tells
you the failure is at the structure axis, fix the schema. When it
tells you the failure is at the environment axis, write more of
yourself into the prompt. When it tells you the failure is
somewhere between, accept the interpretive work, walk both, and
discuss the result with someone who knows the territory. Hold the
ladder loosely: treat it as a vocabulary and an ordering hint,
not as a unique decomposition of every output.

Of all the things this book has tried to give you, this chapter is
the one most likely to change shape under contact with a real user
sitting across a real desk. The other chapters describe how the
machine works. This one describes how to talk about what it gets
wrong, and the talking, in the end, is most of the work.

\section{What Endures and What Ages}
\label{sec:ch13-what-endures}

A reader who returns to this book in three years will find specific
numbers out of date. The observed alignment-tax ranges will be
different. The benchmark scores will be different. The specific
model families will be different. Some of the architectural choices
--- mixture-of-experts routing, specific positional encodings,
particular KV-cache layouts --- will have been superseded.
Re-measure the empirical claims against the current generation;
treat them as calibration data, not as constants.

Three structural observations will not change as quickly, because
they concern the training objectives rather than the specific models
trained under them.

\textbf{The plausibility--truth gap.} As long as language models are
trained on next-token prediction over a corpus, they optimise for
plausibility. Plausibility and truth are correlated where the corpus
is accurate and uncorrelated where the corpus is sparse, wrong, or
silent. Scaling the model and the corpus narrows the practical gap
on the distribution of the training data; it does not close the
conceptual gap, because the objective does not change. The grounding
axis of the error model is the user-visible surface of this
property. It will remain the most consequential axis until the
training objective changes in a fundamental way.

\textbf{The pattern-matcher problem.} A reward model trained on
pairwise preferences is identified only up to an additive constant
per prompt; only ordinal structure is recoverable from the signal.
The policy trained against it learns to produce outputs the reward
model scores highly, which is correlated with but not identical to
what humans actually prefer in situations the reward model has not
seen. The gap between the two --- Goodhart drift, reward hacking,
sycophancy --- is a structural property of training preferences
from comparisons rather than from ground truth. New preference
algorithms narrow the gap; none close it. Knowing this is knowing
why the alignment failure modes of Chapter~9 are not engineering
accidents but structural consequences of the training setup.

\textbf{The form-without-reference problem.} A model trained on
linguistic form has no native mechanism for connecting tokens to
the world the tokens describe. Retrieval, tool use, and grounding
architectures import external reference into the context window at
inference time; they do not solve the underlying problem
architecturally. This is the observation that Harnad's symbol-
grounding problem, Bender and Koller's octopus argument, and a
large body of empirical hallucination research all converge on from
different directions. It is not a claim that the models are useless;
it is a claim about what kind of thing they are and therefore what
kind of failure to expect from them and what kind of engineering
to reach for.

These three are the load-bearing walls of the current architecture.
The error model in this chapter is the user-facing vocabulary for
the same set of facts. A practitioner who holds both --- the
mechanism and the vocabulary --- can navigate the technology as it
changes without having to rebuild intuition from scratch each time
a new model family arrives.
